\subsection{The selector amplifies the penalty it inherits}
Table~\ref{tab:main} gives end-to-end selection. Retrieval into a five-tool
candidate set followed by the 8B selector reaches 54.2\,\% top-1 on the full
1,062-tool catalogue, against a top-5 recall ceiling of 61.1\,\%. The
interesting quantity is how much of that ceiling each subgroup converts.
Conditioned on the gold tool being present in the candidate set, the selector
picks it 95.7\,\% of the time for distinct tools but only 80.3\,\% of the time
for confusable ones. The confusable/distinct gap therefore \emph{grows} from
12.9 points at the retrieval stage to 20.8 points end-to-end: the selector does
not absorb the ambiguity that retrieval passes to it, it adds to it.

This is not an artefact of retrieval. With BM25 candidates the end-to-end gap
is 32.8 points (38.7 vs 71.5), tracking that retriever's larger recall gap;
and removing the candidate bottleneck entirely does not close it. At $N=200$, putting \emph{every} schema in
the prompt raises overall accuracy from 74.0\,\% to 77.8\,\%, but the entire
gain lands on distinct tools (79.7 $\rightarrow$ 85.0); accuracy on confusable
tools is identical to the decimal (58.4 vs 58.4), bought with 4,536 prompt
tokens instead of 211 and 22.8\,s per query instead of 1.1\,s. Twenty-one times the context bought no improvement where the problem is.
We do not claim the distinguishing information is absent from the schema text;
what we measure is that our selector did not exploit it, and that supplying
more of it did not help. Retrieval, usually framed as trading accuracy for
cost, gave up nothing on the confusable subset in our setup.

\subsection{Renaming does not transfer to a standardised catalogue}
PA-Tool-style renaming is the cheapest published remedy for small-model tool
errors~\cite{patool}, and on this catalogue it does not help: 53.2\,\% against
54.2\,\% for the same pipeline with native names, with the confusable subgroup
essentially unchanged (42.9 vs 43.9). The mechanism is visible in the
catalogue itself. Stripping the \texttt{TMFxxx} prefix and snake-casing the
operation makes 293 of 1,062 tools (27.6\,\%) share a name with at least
one other tool, because the standards code \emph{is} the disambiguator between
sibling APIs. Renaming either collides or must re-encode the same information as a domain
prefix, so in our setup pretraining-aligned naming had nothing left to add.

\subsection{Tool knowledge has to be trained in, and it is size-invariant}
In the description-free setting (no schemas in the prompt, roughly 90 tokens) the 8B reasoning model scores 0.2\,\% and the untuned 1.7B scores 0.0\,\%.
Neither can name a TM~Forum operation from pretraining. A LoRA adapter trained
only on template-generated request/tool pairs lifts the 1.7B to 30.8\,\% at
0.30\,s per query; a shorter 300-step run reaches 24.6\,\%, so the recipe is
not saturated at the budget we could afford. The comparison that matters is
not 30.8 against 54.2 but the shape of the curves (Fig.~\ref{fig:retr}):
the router never sees the catalogue, so its accuracy is constant in $N$ by
construction, while the retrieval pipeline falls from 93.8\,\% at $N=25$ to
54.2\,\% at $N=1062$. At full scale the router alone already matches dense
retrieval's top-1 recall (30.8 vs 32.5) without an index.

It is not, however, safe to let it act. Only 43\,\% of its outputs name an
operation that exists, and 22\,\% of queries end in a destructive call that
was not requested ($c_{\mathrm{sev}} = 3.02$, against 1.66 for retrieval).
A description-free router is a cheap \emph{proposer} of candidates and a bad
executor; the accuracy number alone hides this completely.

\subsection{The two signals are complementary}
The hybrid raises accuracy from 54.2\,\% to 62.3\,\% ($+8.1$ points) at the
full catalogue for three extra prompt tokens. The
gain is concentrated where the problem is: confusable accuracy rises from
43.9\,\% to 53.3\,\% ($+9.4$), and candidate recall from 61.1\,\% to
69.5\,\%. The router contributes candidates that retrieval misses rather than
reranking the ones it finds, and routing its proposal through a selector that
can see schemas also repairs its safety profile: the hybrid's $c_{\mathrm{sev}}$
is 1.43 and its dcall 9.8\,\%, both better than dense retrieval alone and far
better than the router acting by itself. A domain-tuned router earns its cost
as a \emph{recall} source at scale, not as a replacement selector; the weaker 300-step adapter already gave $+5.9$.

\subsection{What a wrong tool is, and how to price it}
Of the 192 errors the dense pipeline makes at $N=1062$, 138 name a real tool.
Among those, 77.5\,\% keep the correct verb and 93.5\,\% the correct severity
class, but only 18.1\,\% stay inside the correct API: \emph{right action,
wrong object} is the dominant failure mode. Severity almost never escalates,
so a ``did it escalate'' check reports a clean run; the damage is that 44 of
those 138 errors substitute one \texttt{DELETE} for another on a different
entity. We therefore report \emph{dcall}, the share of queries ending in a
destructive call the request did not ask for: 10.5\,\% for dense top-5.

Scoring this correctly matters. Charging unresolved outputs at read severity, the obvious default, makes the two 0\,\%-accuracy baselines the
safest-looking systems in the table. Once a model emitting
\texttt{TMF630\_\allowbreak removeWorkOrder} is charged for having asked for a deletion
whether or not that operation exists, the ordering inverts: $c_{\mathrm{sev}}$
is 6.78 for the untuned router and 4.51 for the schema-free 8B, against 1.66
for dense retrieval and 1.43 for the hybrid.

\smallskip\noindent\textbf{Arguments fail independently.} With the correct tool supplied, the selector names the right parameters almost always (key precision 94.6\,\%,
recall 98.7\,\%) but gets every value right in only 58.7\,\% of calls
(per-value exact match 77.3\,\%). Composed with selection at the full
catalogue, a complete correct call lands at roughly 32\,\%. Arguments are thus a failure source of comparable size to routing and largely
independent of it, so a single blended ``tool-call accuracy'' would misattribute
half the loss to the router.

\subsection{Abstention, and why accuracy is the wrong headline}
With the gold tool removed and \texttt{NONE} the only correct answer, the
pipeline abstains correctly 29.1\,\% of the time, naming a wrong tool in seven cases out of ten, at the highest severity cost we measured (2.74).
Abstention is worst where it matters: 21.7\,\% with confusers present against
36.7\,\% without. A sibling operation is close enough to look like an answer.
This mirrors the finding that replacing a router model with a retriever preserves
selection quality but cannot provide abstention~\cite{ondevice_routing}, on a
catalogue where close enough may be a \texttt{DELETE}.
